\section{Discussion}\label{discussion}

\subsection{1. What the Privacy-Utility Curve Tells Us About
Deployability}\label{what-the-privacy-utility-curve-tells-us-about-deployability}

The central tension in NGSP is that tighter DP noise (smaller ε)
provides stronger formal guarantees but degrades the proxy text and
therefore the downstream answer quality. The calibration experiment
(\texttt{experiments/calibrate\_epsilon.py}) plots this tradeoff
explicitly.

\textbf{If the curve shows a viable operating point exists (H₁ and H₂
both satisfied at some ε):} ε = 3.0 is likely near the Pareto frontier
for this corpus and model size. Below ε ≈ 1.0, σ \textgreater{} 3.3,
which adds noise large enough relative to the hidden-state magnitude
that the proxy decoder produces semantically incoherent text. Above ε ≈
5.0, the noise is small enough that the inversion attacker may recover
signal from the residual structure.

\textbf{If no viable operating point exists:} This is a meaningful
negative result. It would suggest that Gemma 4 at 2B parameters does not
have sufficient representational capacity to disentangle task intent
from entity-specific content at the hidden-state level --- a finding
that motivates either a larger local model or a different proxy
mechanism (e.g., semantic hashing rather than embedding-space DP).

\textbf{Deployment posture.} Even at a viable operating point, the
system has hard prerequisites for real deployment: - The Gemma model
weights must stay on-premises; cloud loading defeats the purpose. - The
\texttt{entity\_map} must be ephemeral; a persistent map is a
de-anonymization database. - Session budget enforcement must be
tamper-resistant; a compromised session budget counter nullifies the DP
guarantees. - The audit log must be monitored for anomalous canary
appearances; a single canary leak is a deployment failure regardless of
aggregate metrics.

\subsection{2. Where the Composed System
Breaks}\label{where-the-composed-system-breaks}

\subsubsection{Query synthesis path
limitations}\label{query-synthesis-path-limitations}

The \texttt{query\_synthesizer.py} relies on Gemma's ability to extract
task intent without mentioning sensitive entities. For free-form
rewriting tasks (``make this SAE narrative clearer''), the task intent
\emph{is} the content --- there is no separable abstract question to
synthesize. The router should classify these as \texttt{dp\_tolerant} or
\texttt{local\_only}, but routing errors will occasionally send them
through synthesis, producing a synthesized query that either lacks the
context needed for a useful answer or inadvertently preserves
entity-identifying structure.

\subsubsection{DP path proxy decoder
quality}\label{dp-path-proxy-decoder-quality}

The v1 proxy decoder (\texttt{proxy\_decoder.py}) uses the noisy
hidden-state vector only as a soft hint to a Gemma paraphrase call, not
as a hard prefix-tuning signal. This means the DP noise is partially
undone by Gemma's paraphrase, which reconstructs clean prose from the
noisy semantic hint. True prefix-tuning or a dedicated decoder network
trained to map noisy embeddings to text would give stronger guarantees.
The current decoder provides approximate DP properties rather than
strict ones.

\subsubsection{Inversion attacker upper
bound}\label{inversion-attacker-upper-bound}

The DistilBERT inversion attacker in Attack 3 is trained on the same
synthetic distribution as the evaluation set. A real adversary would
likely train on a larger, more varied corpus of (proxy, original) pairs
accumulated over many sessions. The F1 reported here is a lower bound on
what a persistent adversary could achieve; it should be treated as ``can
a naive single-model attacker recover spans?'' rather than ``can any
attacker?''

\subsubsection{Membership inference
scope}\label{membership-inference-scope}

Attack 4 evaluates only the top-3 most-frequent compound codes. Rare
entities (appearing in only a few documents) are harder to evaluate but
also higher-value targets. An adversary specifically interested in a
novel compound code present in only 1--2 documents cannot be evaluated
with the current approach, which requires positive/negative class
balance.

\subsection{3. Honest Deployment
Posture}\label{honest-deployment-posture}

NGSP in its current form is a \textbf{research prototype} appropriate
for: - Evaluating the feasibility of the approach and calibrating the
privacy-utility tradeoff - Benchmarking against simpler baselines (Safe
Harbor only, no-privacy direct send) - Identifying the regime where
formal DP guarantees become practically useful

It is \textbf{not appropriate} for: - Production deployment with real
patient data without independent security review - Regulatory compliance
claims (BAA, HIPAA Safe Harbor certification, etc.) -
Therapeutic-area-specific document types not represented in the
synthetic corpus - Scenarios where the adversary has side-channel access
to model artifacts

\subsection{4. What Is Out of Scope}\label{what-is-out-of-scope}

\textbf{Rare-disease trials.} The synthetic corpus uses
\textasciitilde15 common indication categories. Ultra-rare diseases may
have only a handful of trials globally; quasi-identifier recombination
risk is qualitatively higher, and the Safe Harbor threshold does not
capture it. NGSP's effectiveness in that regime is unknown.

\textbf{Genomics and multi-omics.} The 18 HIPAA identifiers do not cover
genomic quasi-identifiers. Re-identification from genomic proxies
requires a different threat model and is not evaluated here.

\textbf{Real-world non-adversarial testing.} Clinical trial staff
interact with LLM chat interfaces opportunistically, not under
controlled adversarial conditions. Real-world privacy properties (how
often the Safe Harbor stripper misfires, how often Gemma routes
incorrectly, latency impact on user adoption) require a pilot study with
actual users.

\textbf{Adversarial fine-tuning of the cloud model.} If an attacker
controls the cloud LLM (or can observe its training), they could
potentially fine-tune it to elicit entity-specific responses from proxy
inputs. This threat is out of scope for this evaluation.

\subsection{5. Future Work}\label{future-work}

\begin{enumerate}
\def\labelenumi{\arabic{enumi}.}
\item
  \textbf{Stronger proxy decoder.} Train a dedicated text decoder on
  (noisy\_hidden\_state → original\_text) pairs. This would give proper
  DP guarantees for the proxy text rather than approximate ones.
\item
  \textbf{Larger local model.} Repeat the evaluation with Gemma 4 at
  larger parameter counts (9B, 27B) to characterize whether routing
  quality and proxy decoder fidelity improve enough to push the viable
  operating point to tighter ε values.
\item
  \textbf{Adaptive adversary evaluation.} Train the inversion attacker
  with knowledge of the NGSP routing decisions (as if the adversary can
  observe the proxy generation process). This tests the security of the
  mechanism under partial disclosure.
\item
  \textbf{Operational pilot.} Run NGSP as a Chrome extension that
  intercepts clipboard paste events into web-based LLM chat interfaces.
  Measure latency impact, user bypass rate, and false-positive Safe
  Harbor stripping rate on real (non-PHI) clinical writing tasks.
\item
  \textbf{Federated DP accounting.} Extend the session budget to track
  across users (not just sessions), enabling a site-level or
  organization-level privacy budget that degrades gracefully as
  aggregate exposure accumulates.
\end{enumerate}
